\begin{table}[htp]
	\centering
	\renewcommand{\arraystretch}{1.3}
		\begin{tabular}{ | c | c | }
			\hline
			\rowcolor[HTML]{5B9BD5} 
			\textcolor{white}{\textbf{Hiperparámetro}} & \textcolor{white}{\textbf{Espacio de búsqueda}}\\
			\hline
			
			activation (función de activación) & [relu, elu] \\
			\hline
			dropout (tasa de dropout)& [0.0, 0.3] \\
			\hline
			use\_batch\_norm (batch normalization) & [true, false] \\
			\hline
			optimizer & [sgd, adam] \\
			\hline
			lr (learning\_rate) & [0.0001, 0.0005, 0.001, 0.005, 0.01] \\
			\hline
			n\_hidden\_q (número capas ocultas para qgen7) & [2, 5] \\
			\hline
			n\_neurons\_q (neuronas por capa para qgen7) & [64, 256] \\
			\hline
			n\_hidden\_v (número capas ocultas para tensiones) & [1, 2] \\
			\hline
			n\_neurons\_v (neuronas por capa para tensiones) & [64, 256] \\
			\hline
			n\_hidden\_p (número capas ocultas para pérdidas) & [2, 5] \\
			\hline
			n\_neurons\_p (neuronas por capa para pérdidas) & [64, 256] \\
			\hline
			
		\end{tabular} 
		
		\caption{Espacio de búsqueda de hiperparámetros (Red Neuronal MetaGen)}
	\end{table}